% !TEX root = ../main.tex
\chapter{The Layered Dependency DAG}
\label{ch:layered_dag}

\partref{part:hardware} described the substrate on which \Lucid\ runs.
\partref{part:system} describes \Lucid\ itself, beginning here with
the most consequential single piece of architecture in the engine:
the layered dependency DAG that organises every C++ source file and
every Python module in the codebase.

The DAG is not a documentation artefact. It is enforced by the build
system, by code-review policy (the \emph{Hard Rules} of
\appref{app:hard_rules}), and by a class of automated checks that
run on every pull request. The chapter that the reader is now reading
exists because the DAG is the thing that, once internalised, makes
the rest of the engine code easy to navigate. A new contributor who
understands the seven layers and their permitted edges can locate
any operation, any class, any allocation in roughly the time it takes
to skim a directory listing. A contributor who has not internalised
it spends hours hunting for definitions that turn out to live two
layers below where intuition suggested.

We proceed bottom up. \secref{sec:layered_dag:layers} describes the
seven layers and their contents. \secref{sec:layered_dag:rules}
states the edge rules that the DAG enforces.
\secref{sec:layered_dag:enforcement} describes how the rules are
enforced --- through CMake target boundaries, through code review,
and through a small set of static checks.
\secref{sec:layered_dag:python_mirror} treats the Python side, which
mirrors the C++ DAG with a smaller set of layers.
\secref{sec:layered_dag:crosscutting} catalogues the cross-cutting
concerns (allocator, error, dispatch) that touch multiple layers
without violating the DAG. \secref{sec:layered_dag:adding} describes
the procedure for adding a new layer or a new file within an existing
layer.

\section{Why a Layered Architecture}
\label{sec:layered_dag:why}

The engine is roughly \(45{,}000\) lines of C++ across more than
two hundred translation units, plus a comparable volume of Python
code that mirrors it. At that scale, an architecture without
explicit structural rules tends to collapse into an undirected
dependency graph, with each module reaching into every other module
that has a useful symbol. Two years into the project, every file
includes every other file; build times explode; refactoring becomes
impossible; new contributors cannot tell where a feature should
live.

The cure is a directed acyclic graph of dependencies, with a small
number of layers, and with an enforced rule that edges go only
downward. This is not a novel pattern --- it is the standard
operating-system-kernel layout, the standard compiler-construction
pattern, the standard database-system layout --- but it is unusual
in ML frameworks, which tend to grow organically and tolerate
arbitrary internal cycles.

The benefits, concretely, for \Lucid:

\begin{enumerate}
  \item \textbf{Compile-time isolation.} A change in the highest
    layer (\code{bindings}) recompiles only that layer; a change in
    the lowest layer (\code{core}) recompiles everything. The CMake
    target graph mirrors the DAG, so this isolation is automatic.
  \item \textbf{Testability.} Each layer can be tested against
    mocks of the layer immediately below. \Lucid's C++ unit tests
    follow this pattern: \code{test/core/} tests \code{core} in
    isolation, \code{test/backend/} tests \code{backend} against
    in-memory mocks of \code{core}, and so on.
  \item \textbf{Mental model.} A contributor reading source code
    can rely on the rule that any symbol referenced is defined in
    the current layer or strictly below. This eliminates the
    ``where could this possibly come from'' search that pervades
    cyclic codebases.
  \item \textbf{Refactoring discipline.} Changes to a layer's
    public API force consumers in higher layers to adapt, in a
    direction that the build system actually checks. This is the
    inverse of the situation in cyclic codebases, where a
    refactoring must be performed simultaneously across every
    affected file.
\end{enumerate}

The cost is up-front structural design and ongoing discipline:
contributors must occasionally accept that the right place for a
new piece of code is not the most convenient one. The \emph{Hard
Rules} (\appref{app:hard_rules}) encode the discipline.

\section{The Seven Layers}
\label{sec:layered_dag:layers}

The C++ engine is organised into seven layers, each occupying a
directory under \code{lucid/\_C/}. We describe them from the bottom
up. The corresponding Python tree, which is shallower, is treated
separately in \secref{sec:layered_dag:python_mirror}.

\figref{fig:layered_dag:overview} gives the full stack at a glance;
the subsequent subsections expand each layer.

\begin{figure}[t]
\centering
\begin{adjustbox}{max width=\textwidth, center}
\begin{tikzpicture}[
  layer/.style={draw, rounded corners=2pt, minimum width=12cm,
                minimum height=22pt, font=\small\sffamily,
                inner sep=4pt, align=left},
  l0/.style={layer, fill=lucid.note.bg},
  l1/.style={layer, fill=lucid.ok.bg},
  l2/.style={layer, fill=lucid.info.bg},
  l3/.style={layer, fill=lucid.code.bg},
  arrow/.style={->, >=stealth, thick, color=lucid.accent.dark,
                shorten >=2pt, shorten <=2pt},
]
  \node[l0] (core) at (0, 0)
    {\textbf{L0 core} \quad Dtype, Device, Shape, Allocator, Error, GradMode, AmpPolicy};
  \node[l0] (tensor) at (0, 1.0)
    {\textbf{L1 tensor} \quad TensorImpl, Storage, AutogradMeta};
  \node[l1] (backend) at (0, 2.0)
    {\textbf{L2 backend} \quad IBackend, Dispatcher, AccelerateBackend, MlxBackend};
  \node[l1] (kernel) at (0, 3.0)
    {\textbf{L3 kernel} \quad UnaryKernel, BinaryKernel, ReduceKernel, primitives/};
  \node[l2] (autograd) at (0, 4.0)
    {\textbf{L4 autograd} \quad Engine, Node, AccumulateGrad, CustomFunction};
  \node[l2] (ops) at (0, 5.0)
    {\textbf{L5 ops / nn / optim / compile} \quad $\sim$260 ops, NN kernels, optimisers, MpsBuilder};
  \node[l3] (bindings) at (0, 6.0)
    {\textbf{L6 bindings} \quad pybind11 module: \texttt{lucid.\_C.engine}};

  \draw[arrow] (tensor) -- (core);
  \draw[arrow] (backend) -- (tensor);
  \draw[arrow] (kernel) -- (backend);
  \draw[arrow] (autograd) -- (kernel);
  \draw[arrow] (ops) -- (autograd);
  \draw[arrow] (bindings) -- (ops);
\end{tikzpicture}
\end{adjustbox}
\caption[The Lucid engine DAG.]{The seven-layer dependency DAG of
the \Lucid\ C++ engine. Solid arrows indicate direct dependency
(layer above includes from layer below). No back-edges (arrows from
a lower layer to a higher one) are permitted; the absence is
enforced by the CMake target graph and verified by
\texttt{tools/check\_dag.py}. Cross-cutting concerns (allocator,
error, dispatch) are defined in \texttt{core} (L0) and consumed
from every layer above.}
\label{fig:layered_dag:overview}
\end{figure}

The DAG is not a documentation idealisation. It corresponds
directly to the directory layout under \code{lucid/\_C/}, to the
CMake target structure, and to the build's link line.
\tabref{tab:layered_dag:files} summarises the per-layer files and
their approximate volumes.

\begin{table}[t]
\centering
\small
\begin{tabular}{llrl}
\toprule
Layer & Directory & SLOC & Key types / files \\
\midrule
L0 core      & \texttt{lucid/\_C/core/}     & $\sim$1\,800  & Dtype, Device, Shape, Allocator, LucidError \\
L1 tensor    & \texttt{lucid/\_C/tensor/}   & $\sim$1\,200  & TensorImpl, Storage, AutogradMeta \\
L2 backend   & \texttt{lucid/\_C/backend/}  & $\sim$6\,500  & IBackend, Dispatcher, two backends \\
L3 kernel    & \texttt{lucid/\_C/kernel/}   & $\sim$3\,400  & Kernel templates + primitives \\
L4 autograd  & \texttt{lucid/\_C/autograd/} & $\sim$4\,900  & Engine, Node, CustomFunction \\
L5 ops       & \texttt{lucid/\_C/ops/}      & $\sim$11\,000 & $\sim$260 op implementations \\
L5 nn        & \texttt{lucid/\_C/nn/}       & $\sim$4\,800  & Conv, BatchNorm, LSTM, Attention \\
L5 optim     & \texttt{lucid/\_C/optim/}    & $\sim$1\,700  & Optimizer family + LR schedulers \\
L5 compile   & \texttt{lucid/\_C/compile/}  & $\sim$5\,200  & Tracer, MpsBuilder, OpEmitters \\
L6 bindings  & \texttt{lucid/\_C/bindings/} & $\sim$3\,500  & pybind11 wrappers \\
\midrule
Total        &                              & $\sim$45\,000 & \\
\bottomrule
\end{tabular}
\caption[Engine SLOC by layer.]{Approximate source-lines-of-code
per engine layer, as of \Lucid\ 3.5. Layer~5 dominates the line
count because it hosts the \(\sim 260\) operation implementations
and the model-zoo helpers.}
\label{tab:layered_dag:files}
\end{table}

\subsection{Layer 0: core}
\label{ssec:layered_dag:core}

\code{lucid/\_C/core/} contains the foundational types on which
every other layer depends. It has no dependencies on any other
\Lucid\ code; only on the C++ standard library and on Apple's
\Accelerate\ (for low-level math helpers). The contents are
deliberately minimal:

\begin{itemize}
  \item \code{Dtype.h}: the dtype enum and conversion helpers.
  \item \code{Device.h}: the device enum (CPU / Metal) and a small
    set of capability queries.
  \item \code{Shape.h}, \code{Shape.cpp}: the shape vector type
    with broadcast and validation operations.
  \item \code{Allocator.h}, \code{Allocator.cpp}: the size-class
    allocator pool.
  \item \code{MemoryStats.h}: per-device allocation statistics.
  \item \code{Error.h}, \code{ErrorBuilder.cpp}: the
    \code{LucidError} exception class and the fluent builder used
    by every error site (\appref{app:pillars} for the production
    pillar P1 that motivates it).
  \item \code{Determinism.h}, \code{Determinism.cpp}: the
    deterministic-kernel-selection flag.
  \item \code{GradMode.h}, \code{GradMode.cpp}: the
    \code{no\_grad} / \code{enable\_grad} state.
  \item \code{AmpPolicy.h}, \code{AmpPolicy.cpp}: the mixed
    precision dtype-casting rules.
  \item \code{Generator.h}, \code{Generator.cpp}: the random
    number generator state.
\end{itemize}

The discipline: nothing in \code{core} may include from any other
\code{lucid/\_C/} directory. Violating this rule introduces a
dependency cycle, since every other layer depends on \code{core}.

\subsection{Layer 1: tensor}
\label{ssec:layered_dag:tensor}

\code{lucid/\_C/tensor/} contains the tensor model itself:

\begin{itemize}
  \item \code{TensorMeta.h}: the \code{TensorImpl} class --- shape,
    stride, dtype, device, storage offset, and a handle to the
    backing storage.
  \item \code{AutogradMeta.h}: the per-tensor autograd metadata
    (\code{requires\_grad}, \code{grad\_fn}, backward hooks).
  \item \code{Storage.h}: the reference-counted storage object
    that wraps a backing buffer (allocated through \code{core}'s
    \code{Allocator}).
\end{itemize}

\code{tensor} depends on \code{core}. It does not depend on any
backend (since \code{TensorImpl} is backend-agnostic) and does not
know how operations are dispatched (that is the next layer's job).
The chapter that fills in this layer in depth is
\chref{ch:tensor_model}.

\subsection{Layer 2: backend}
\label{ssec:layered_dag:backend}

\code{lucid/\_C/backend/} contains the abstract backend interface
and the two concrete implementations:

\begin{itemize}
  \item \code{IBackend.h}: the pure-virtual interface that every
    backend implements: \code{allocate}, \code{deallocate},
    elementwise op categories, reductions, matmul, convolution,
    and so on.
  \item \code{Dispatcher.h}, \code{Dispatcher.cpp}: the router
    that selects a backend based on a tensor's device.
  \item \code{cpu/AccelerateBackend.\{h,cpp\}} and several
    helper files: the \Accelerate-backed CPU backend
    (\chref{ch:accelerate}).
  \item \code{gpu/MlxBackend.\{h,cpp\}}: the \MLX-backed GPU
    backend (\chref{ch:mlx_runtime}).
  \item \code{gpu/SharedStorage.\{h,cpp\}}: the zero-copy
    CPU/GPU bridge (\chref{ch:unified_memory}).
\end{itemize}

\code{backend} depends on \code{core} and \code{tensor}. It is the
first layer that touches external libraries (\Accelerate, \MLX) at
build time. The chapter that fills in this layer is
\chref{ch:backend_dispatch}.

\subsection{Layer 3: kernel}
\label{ssec:layered_dag:kernel}

\code{lucid/\_C/kernel/} contains the kernel abstraction --- the
templates that take a backend-agnostic operation and dispatch it
through the backend's primitives:

\begin{itemize}
  \item \code{IKernel.h}: the abstract base.
  \item \code{UnaryKernel.h}, \code{BinaryKernel.h},
    \code{ReduceKernel.h}, \code{NaryKernel.h},
    \code{VariadicKernel.h}: the templates per arity class.
  \item \code{Contig.h}: contiguity inference and conversion.
  \item \code{primitives/}: data-dependent primitives that cannot
    be expressed as plain elementwise or reduction operations
    (Gather, Scatter, Im2Col, Pad, Sort, NonZero, Unique,
    SearchSorted, and others).
\end{itemize}

\code{kernel} depends on \code{core}, \code{tensor}, and
\code{backend}. The chapter is \chref{ch:kernel_abstraction}.

\subsection{Layer 4: autograd}
\label{ssec:layered_dag:autograd}

\code{lucid/\_C/autograd/} contains the differentiation engine:

\begin{itemize}
  \item \code{Node.h}, \code{Node.cpp}: the base class for
    backward nodes.
  \item \code{AutogradNode.h}: the templated convenience class
    used by most ops.
  \item \code{Engine.h}, \code{Engine.cpp}: the reverse-mode
    backward walker (topological sort + per-node \code{apply}).
  \item \code{AccumulateGrad.\{h,cpp\}}: the leaf node that holds
    accumulated gradients on parameters.
  \item \code{CustomFunction.\{h,cpp\}}, \code{FuncOp.h}: the
    user-extensible custom-function machinery.
  \item \code{FusionPass.\{h,cpp\}}: a forward-only fusion pass
    that combines compatible adjacent ops (e.g.\ Linear+ReLU,
    SDPA).
\end{itemize}

\code{autograd} depends on \code{core}, \code{tensor},
\code{backend}, and \code{kernel}. It is responsible for
constructing the backward graph as ops execute and for traversing
it on \code{backward()}. \chref{ch:autograd_engine} fills in the
details.

\subsection{Layer 5: ops, nn, optim, linalg, fft, einops, complex,
compile}
\label{ssec:layered_dag:ops}

The fifth layer is broad. It contains the actual operation
implementations, grouped into several sibling directories:

\begin{itemize}
  \item \code{lucid/\_C/ops/}: the elementwise (ufunc, bfunc),
    generic (gfunc), composite, linalg, fft, einops, and complex
    operations. Each operation is a single C++ file that takes
    backend-agnostic \code{TensorImpl} inputs, dispatches through
    \code{backend} or constructs through \code{kernel}, and
    returns one or more outputs. There are roughly 260 such
    operations.
  \item \code{lucid/\_C/nn/}: neural-network-specific kernels
    (Conv, BatchNorm, GroupNorm, LayerNorm, Embedding, LSTM,
    MultiheadAttention, Interpolate, AdaptivePool, Dropout). These
    differ from \code{ops/} in that they often manage internal
    running statistics or expose neural-network-specific APIs.
  \item \code{lucid/\_C/optim/}: the optimiser implementations
    (Optimizer base class, SGD, Adam family, AdaGrad family,
    RMSprop, LBFGS, and a number of LR schedulers).
  \item \code{lucid/\_C/compile/}: the graph-capture compiler ---
    Tracer, TraceIR, MpsBuilder, ExecutableCache, OpEmitters
    (\chref{ch:compile_design}).
\end{itemize}

These siblings all depend on \code{core} through \code{autograd}.
They do not depend on each other; an \code{ops/} file does not
include from \code{nn/} or \code{optim/}, and vice versa. The
chapters that fill in the contents of layer 5 are
\chref{ch:op_registry} through \chref{ch:lr_schedulers} for
operations and optimisers, and \partref{part:graph_compile} for
the compile path.

\subsection{Layer 6: bindings}
\label{ssec:layered_dag:bindings}

\code{lucid/\_C/bindings/} contains the pybind11 binding code that
exposes the C++ engine to Python. Each binding file wraps a
related set of operations:

\begin{itemize}
  \item \code{bind.cpp}: the module bootstrap.
  \item \code{bind\_tensor.cpp}: the \code{Tensor} class.
  \item \code{bind\_bfunc.cpp}, \code{bind\_ufunc.cpp},
    \code{bind\_gfunc.cpp}, \code{bind\_composite.cpp}: per-op
    bindings.
  \item \code{bind\_autograd.cpp}, \code{bind\_nn.cpp},
    \code{bind\_optim.cpp}, \code{bind\_linalg.cpp},
    \code{bind\_fft.cpp}, \code{bind\_einops.cpp},
    \code{bind\_complex.cpp}, \code{bind\_amp.cpp},
    \code{bind\_profiler.cpp}, \code{bind\_random.cpp},
    \code{bind\_compile.cpp}: subsystem-level bindings.
  \item \code{bind\_errors.cpp}: the exception translation that
    maps \code{LucidError} to a Python exception.
  \item \code{bind\_op\_registry.cpp}: a binding for the op
    registry that introspection tools can read.
\end{itemize}

\code{bindings} depends on every layer below it. It is the only
layer that links against pybind11 itself; the rest of the engine
is pure C++. This isolation matters: \Lucid's engine could in
principle be linked into a non-Python program by replacing the
bindings layer with a C API, and we have kept the option open by
ensuring no pybind11 types leak into lower layers.

\chref{ch:python_cpp_binding} fills in the binding details.

\section{Dependency Rules}
\label{sec:layered_dag:rules}

The dependency rules between layers are simple and absolute.

\subsection{Acyclicity}
\label{ssec:layered_dag:acyclic}

The dependency graph between layers is acyclic. There is no edge
from a lower-numbered layer back up to a higher-numbered one. The
ordering is total: \code{core} $\prec$ \code{tensor} $\prec$
\code{backend} $\prec$ \code{kernel} $\prec$ \code{autograd}
$\prec$ \code{ops}/\code{nn}/\code{optim}/\code{compile} $\prec$
\code{bindings}.

The ordering implies, for any pair of layers $A$ and $B$ with
$A \prec B$:

\begin{itemize}
  \item $B$ may include headers from $A$.
  \item $B$ may call functions defined in $A$.
  \item $A$ may not include headers from $B$.
  \item $A$ may not name any type defined in $B$.
  \item $A$ may not call any function defined in $B$.
\end{itemize}

\subsection{Edge directness}
\label{ssec:layered_dag:directness}

Edges are typically direct: layer $i+1$ depends on layer $i$, not
on layer $i-2$. But edges may legitimately skip layers when the
intermediate layer does not own the abstraction in question. For
example, \code{ops/Reduction.cpp} (layer 5) depends directly on
\code{core/Dtype.h} (layer 0) and on \code{kernel/ReduceKernel.h}
(layer 3) but not on \code{tensor/} or \code{backend/} headers,
because the reduction is expressed in terms of kernel templates
that themselves handle tensor and backend concerns.

Skipping layers is fine. Reversing direction is not.

\subsection{Sibling independence within layer 5}
\label{ssec:layered_dag:siblings}

Within layer 5, the four sibling directories (\code{ops/},
\code{nn/}, \code{optim/}, \code{compile/}) do not depend on each
other. This is not a logical necessity --- one could imagine
\code{nn/} reusing some \code{ops/} composition --- but it is an
empirical discipline that prevents the layer from collapsing into
a tangle. When a \code{nn/} kernel needs a primitive that
\code{ops/} provides, it depends on the layer-4 \code{autograd}
representation of the primitive, not on \code{ops/} directly.

\code{compile/} is a partial exception: it must enumerate every op
its emitter supports, which means it has read-only knowledge of
the names and signatures in \code{ops/} and \code{nn/}. We
implement this by having \code{compile/OpEmitters/} hold a table
of emitter functions keyed by op name; the table is populated at
build time by code generation from a single source of truth in
\code{compile/OpRegistry.h}, not by header includes.

\section{Enforcement}
\label{sec:layered_dag:enforcement}

The DAG would be a documentation artefact if it were not enforced.
\Lucid\ enforces it through three independent mechanisms.

\subsection{CMake target boundaries}
\label{ssec:layered_dag:cmake}

Each layer is a separate CMake target. The CMakeLists structure
under \code{lucid/\_C/} declares one library per layer, with
explicit \code{target\_link\_libraries} edges that mirror the
permitted dependencies. Trying to use a symbol from a layer that
is not in the current target's link line produces a link error,
which is caught at build time rather than at code review.

This is the strongest enforcement mechanism. A contributor who
attempts to violate the DAG sees a build failure within seconds of
saving the file.

\subsection{Code review policy}
\label{ssec:layered_dag:review}

The Hard Rules of \appref{app:hard_rules} include an explicit
rule (H-rule on layer dependency) that any pull request which
adds an edge in the wrong direction is rejected. Reviewers are
expected to verify, for each new include in a header, that the
include's directory is at or below the current file's layer.

In practice this rarely needs invocation, because the CMake
boundary catches the violation first. But the rule exists for the
edge case in which a violation passes CMake (for example, a
contributor adds the offending edge to the CMake target as part of
the same PR), and it gives the reviewer explicit grounds to
reject.

\subsection{The \texttt{tools/check\_dag.py} script}
\label{ssec:layered_dag:check_dag}

A small static-analysis script in \code{tools/check\_dag.py}
parses the include statements in every C++ header and source file
under \code{lucid/\_C/} and verifies that every include's target
directory is at or below the source file's layer. The script is
run as part of the pre-commit hook and as part of the CI check
list; a violation fails the build.

The script is simple (under two hundred lines of Python) and uses
only the standard library. It is the fallback in case a CMake
restructuring inadvertently relaxes the target boundary.

\section{The Python Mirror}
\label{sec:layered_dag:python_mirror}

The Python side of \Lucid\ mirrors the C++ DAG with a smaller set
of layers. The structure is:

\begin{itemize}
  \item \textbf{Foundation:} \code{lucid/\_types.py},
    \code{lucid/\_dtype.py}, \code{lucid/\_device.py},
    \code{lucid/\_globals.py}. Pure Python, no engine dependency.
  \item \textbf{Engine binding:} \code{lucid/\_C/} (the compiled
    extension module produced by the C++ layers). Imported under
    the \code{\_C\_<name>} alias convention required by Hard Rule
    H2.
  \item \textbf{Dispatch:} \code{lucid/\_dispatch.py}, the single
    boundary that wraps and unwraps between Python \code{Tensor}
    and engine-side \code{TensorImpl}.
  \item \textbf{Tensor:} \code{lucid/\_tensor/} (the \code{Tensor}
    class, methods, dunders, indexing, conversion, repr).
  \item \textbf{Factories:} \code{lucid/\_factories/} (tensor
    creation: \code{zeros}, \code{ones}, \code{tensor},
    \code{from\_numpy}, \code{from\_dlpack}, random).
  \item \textbf{Op wrappers:} \code{lucid/\_ops/} (the registry
    of free functions and Tensor methods exposed at the top level).
  \item \textbf{Subsystem modules:} \code{lucid/autograd/},
    \code{lucid/nn/}, \code{lucid/optim/}, \code{lucid/linalg/},
    \code{lucid/fft/}, \code{lucid/signal/},
    \code{lucid/special/}, \code{lucid/distributions/},
    \code{lucid/func/}, \code{lucid/einops/}, \code{lucid/amp/},
    \code{lucid/profiler/}, \code{lucid/serialization/},
    \code{lucid/compile/}, \code{lucid/models/}.
  \item \textbf{Top-level package:} \code{lucid/\_\_init\_\_.py},
    the lazy-loading dispatch surface
    (\chref{ch:lazy_imports}).
\end{itemize}

The Python rules are looser than the C++ rules, because Python
permits cycles that the C++ build would reject, but the convention
is the same: factories depend on engine binding, tensor depends on
factories + engine binding, subsystem modules depend on
tensor/factories/engine, and the top-level package depends on
everything.

The convention is enforced by code review rather than by automation,
because a comprehensive Python import-graph checker would have to
parse Python's dynamic import semantics, which is more work than
the benefit justifies.

\section{Cross-cutting Concerns}
\label{sec:layered_dag:crosscutting}

A handful of concerns touch every layer without violating the DAG.
The pattern is the same in each case: the concern is defined in
\code{core} and consumed by every layer above. Three are worth
naming explicitly.

\subsection{The allocator}
\label{ssec:layered_dag:allocator_xc}

\code{core/Allocator} is the size-class pool that every backend's
memory allocations route through. Layer 2 (\code{backend}) wraps
the underlying provider's allocator (\MLX's or \Accelerate's) and
funnels requests through the pool. Layer 5 (\code{ops}) and beyond
do not allocate directly; they always go through the backend.

The chapter is \chref{ch:allocator}.

\subsection{Error handling}
\label{ssec:layered_dag:errors_xc}

\code{core/Error} defines \code{LucidError}, the single exception
type that propagates from any layer back to Python. Layer 6
(\code{bindings}) installs the pybind11 translator that maps
\code{LucidError} to the Python \code{LucidError} class.
Intermediate layers may catch and rethrow with additional context;
the chapter is \chref{ch:error_handling}.

\subsection{Dispatch}
\label{ssec:layered_dag:dispatch_xc}

The dispatcher (layer 2, \code{backend/Dispatcher.h}) is the
single function that routes per-op requests to the correct backend.
Above it, \code{ops/} files do not call backends directly; they
call the dispatcher with the op kind and the input tensors, and
the dispatcher selects the implementation. This indirection is
what allows the same op file to work for both CPU and GPU without
conditional code.

The chapter is \chref{ch:backend_dispatch}.

\section{Adding a Layer or a File}
\label{sec:layered_dag:adding}

The protocol for adding a new file within an existing layer is
short:

\begin{enumerate}
  \item Identify the layer that the new file belongs to. The
    criterion is: what does the file depend on? If it depends only
    on \code{core}, it lives in \code{tensor} or higher; if it
    depends on \code{tensor}, in \code{backend} or higher; and so
    on.
  \item Add the file under the appropriate directory.
  \item Update the CMakeLists in that directory to include the
    new file in the layer's library target.
  \item Add the file's tests to the corresponding
    \code{test/<layer>/} directory.
  \item Confirm the build still passes and the check\_dag script
    accepts the new file.
\end{enumerate}

Adding a new layer is much rarer and requires explicit design
discussion. The reason: a new layer either fits between two
existing ones (in which case it adds a hop to every dependency
chain through that point) or sits at the top above
\code{bindings} (in which case it suggests the engine has grown a
new audience). Both situations have happened exactly once each in
\Lucid's history. The first was the addition of \code{kernel} as a
layer above \code{backend}, which absorbed the kernel-template
logic that had previously lived in individual \code{ops/} files
and was being duplicated. The second was the introduction of
\code{compile/} as a sibling of \code{ops/}, \code{nn/}, and
\code{optim/}, which was added to host the graph-capture compiler
without polluting the eager-mode hierarchy.

\section{Summary and Forward References}
\label{sec:layered_dag:summary}

The layered DAG is the structural invariant that makes the \Lucid\
engine tractable. Seven layers, edges only from higher to lower,
enforced by the CMake target graph plus a static analysis script
plus code-review discipline. Every other architectural decision in
the engine fits within the DAG; every cross-cutting concern is
located in \code{core} and consumed by every layer above; the
Python side mirrors the C++ side with a smaller and looser set of
rules.

The chapters that follow expand on individual layers.
\chref{ch:tensor_model} describes layer 1 in detail.
\chref{ch:allocator} treats the cross-cutting allocator.
\chref{ch:backend_dispatch} treats layer 2 and the dispatcher.
\chref{ch:bridge_boundaries} catalogues the six places where the
engine touches the outside world.
\chref{ch:python_cpp_binding} returns to layer 6 with a focus on
the pybind11 contract. \chref{ch:lazy_imports} closes
\partref{part:system} with the Python-side lazy-import surface
that the top-level package uses to make \Lucid\ importable in
under one hundred milliseconds despite the size of its subsystem
tree.

\begin{lucidnote}
The DAG is invisible at the user surface --- the user imports
\code{lucid}, constructs a model, calls \code{forward}, and never
sees layer boundaries. The DAG matters only to contributors and to
the build system. But its presence is the reason \Lucid\ has not
collapsed into the tangle of mutual includes that ML frameworks of
comparable scope sometimes become.
\end{lucidnote}
